% 第12回　最終プレゼン発表会（記入例・模範解答）
\session{12}{2}{2027年1月13日（水）}{最終プレゼン発表会}{}{yes}

\goals{
  \item 企業の担当者に向けて、店舗設計の提案をわかりやすく伝える
  \item AIをどこでどう使い、何を採用し何を捨てたか、自分たちの思考がどこに反映されているかを報告する
}

\work[50mm]{発表前のチェック}{当日までに}
\begin{achecks}
  \item スライドは指定の方法で提出し、当日使う機材で表示を確認した
  \item 発表時間内に収まる
  \item 参考にしたデータや事例の出典が示してある
  \item 生成画像には「AIによる生成イメージ」と明記してある
  \item 「AI活用と自分たちの思考」のスライドがある
  \item 質問されそうな点への答えを用意した
\end{achecks}

\work[60mm]{準備}{想定質問と答え}
\begin{wstabh}{colspec={X[l,m] X[l,m] Q[l,wd=16mm,m]},row{2-Z}={ht=13mm}}
  想定される質問 & 答え（根拠） & 答える人 \\
  \af{話しかけない接客で、売上は下がらないのか} & \af{目的は来店のハードルを下げること。相談したい客には呼ばれたときに集中して対応できる。まず1か月試して記録する提案} & \af{自分} \\
  \af{改装の費用はどのくらいかかるのか} & \af{具体的な金額は示せない。費用の小さい施策（カード・サイン・配置換え）から始める3段階の案} & \af{Ｄ} \\
  \af{既存の客（整備・年配の客）は離れないか} & \af{整備の静かな待合を残す。既存客のペルソナでも動線を検証した} & \af{Ｃ} \\
  \af{キッズスペースの安全は} & \af{展示車の固定、柵の設置、車の出し入れは閉店後。休日は見守り担当を置く} & \af{Ｂ} \\
\end{wstabh}

\work[40mm]{確認}{発表で見られること（自己評価 ○△×）}
\begin{wstabh}{colspec={X[l,m] Q[c,wd=18mm,m] Q[c,wd=18mm,m]},row{2-Z}={ht=7mm}}
  観点 & 発表前 & 発表後 \\
  企業の課題を正しく捉えているか & \ans{○} & \ans{○} \\
  ペルソナとコンセプトに一貫性があるか & \ans{○} & \ans{○} \\
  空間と動線の提案に根拠があるか & \ans{△} & \ans{△} \\
  実現性と費用を考えているか & \ans{△} & \ans{×} \\
  AIの活用が透明で、自分たちの判断を説明できるか & \ans{○} & \ans{○} \\
\end{wstabh}

\work[70mm]{当日}{他グループの発表メモ}
\begin{wstabh}{colspec={Q[l,wd=20mm,m] X[l,m] X[l,m]},row{2-Z}={ht=15mm}}
  グループ & 印象に残った点 & 質問・自分たちとの違い \\
  \af{Ａ} & \af{夜だけ開くガレージ。時間帯で客層を分ける} & \af{スタッフの勤務の工夫まで説明していた。自分たちは運用の説明が弱かった} \\
  \af{Ｂ} & \af{オンライン予約と来店を組み合わせた「答え合わせ」の店} & \af{前半の自分の問いと近い。来店の目的に注目している点が共通} \\
  \af{Ｄ} & \af{整備待ちの時間を体験に変える} & \af{既存客から始める点が現実的。費用の見積もりの考え方が具体的} \\
  & & \\
\end{wstabh}

\work{講評}{企業担当者からの講評・質疑で受けた指摘}
\alines{4}{・ペルソナの言葉から空間までが一貫していて、わかりやすい。\par
・「声をかけない」は現場では勇気がいる。スタッフの評価の仕組みも合わせて考えてほしい。\par
・費用の説明が弱い。金額でなくても、段階ごとの規模感は示せたはず。\par
・AIを使った場面と捨てた理由が明確で、判断の跡が見えた。}

\work{振り返り}{発表を終えて（うまくいったこと／次に活かすこと）}
\atwoboxes{うまくいったこと}{次に活かすこと}{30mm}
{ペルソナの言葉で始め、最後まで同じ2人の視点で説明できた。想定質問の準備が当たり、安全の質問に落ち着いて答えられた。}
{費用や運用（スタッフの評価）のような、企業側が判断に使う情報を早めに確認する。数字がなくても規模感を示す方法を考える。}

\work[30mm]{提出物（後半）}{店舗設計提案}
\begin{achecks}
  \item スライド（必要に応じて補足資料）
  \item AI活用記録を添付した
\end{achecks}
\afillin{書式・締切（授業内で案内）}{（授業内で案内した書式・締切を記入）}

\begin{point}
  \item 当日の評価は、授業計画に示した5つの観点（課題の把握、ペルソナとコンセプトの一貫性、空間と動線の根拠、実現性と費用、AI活用の透明性）で行う。発表後の自己評価と照らし合わせて、講評の際に伝える。
  \item 想定質問は、企業側の観点（費用、既存客、安全、運用）から準備できているかを見る。
  \item 他グループの発表メモは、自分たちとの違いまで書けていると、第13回の振り返りに生きる。
  \item 企業担当者の講評は、その場でメモさせ、提出物の改訂にどう反映したかを確認する。
\end{point}
